\documentclass{article}

\usepackage[utf8]{inputenc}
\usepackage[T1]{fontenc}
\usepackage{amsmath,amssymb,amsthm}
\usepackage{mathtools}
\usepackage{hyperref}
\usepackage{booktabs}

\makeatletter
\newcommand{\citep}[2][]{%
  \begingroup
  \def\@tmpkey{#2}%
  \ifx\@tmpkey\@empty\else
    (\ifx\@empty#1\@empty\else #1, \fi
     \@nameuse{bibkey@\@tmpkey})%
  \fi
  \endgroup
}
\expandafter\def\csname bibkey@lasser2026micro\endcsname{Microfoundation}
\expandafter\def\csname bibkey@lasser2026exo\endcsname{Exogenous~Verification}
\expandafter\def\csname bibkey@lasser2026phase\endcsname{Phase~Redundancy}
\makeatother

\newtheorem{definition}{Definition}
\newtheorem{lemma}{Lemma}
\newtheorem{theorem}{Theorem}
\newtheorem{remark}{Remark}
\newtheorem{assumption}{Assumption}
\newtheorem{proposition}{Proposition}

\newcommand{\Bclip}{B_{\mathrm{clip}}}
\newcommand{\Kch}{K_{\mathrm{ch}}}
\newcommand{\Aadv}{\mathcal{A}_{\mathrm{adv}}}
\newcommand{\Pproxy}{P^{\mathrm{proxy}}}
\newcommand{\Ttruth}{T^{\mathrm{truth}}}
\newcommand{\epsnonresCoev}{\varepsilon_{\mathrm{coev}}^{\mathrm{nonres}}}
\newcommand{\Nev}{N_{\mathrm{ev}}}
\newcommand{\rhogap}{\rho_{\mathrm{gap}}}

\title{C7 closed-form derivation: bounding $\bar{M}$ from primitive operational parameters}
\author{Paper 10 working draft for codex review}
\date{}

\begin{document}

\maketitle

\section{Goal of this draft}

Currently Paper 10's (C7) bounded co-evolution
(Assumption~\ref{ass:bounded_coev}) is asserted as an explicit
deployment-class assumption: there exists a constant $\bar{M}$
depending on deployment-class parameters but independent of $|P|$
such that the per-channel coupling magnitude
$M(\mathrm{deployment}) \leq \bar{M}$.  Audit~4
(\S\ref{sec:tooling_substrate_audits}) gives an empirical procedure
for measuring $\bar{M}$ but does not derive it from primitive
operational parameters.

This draft attempts to close the gap by deriving $\bar{M}$ in
closed form from already-named deployment-class quantities ---
specifically, from (C5.HOEFF) clip radius $\Bclip$, (C5.MULT)
channel partition structure, and a verification-protocol-derived
\emph{shared-action probability bound} $Q^*$.  If the derivation
holds, (C7) graduates from assumption to corollary, and the loose
thread C7 represents in the deployment-claim hypothesis set
disappears.

The derivation does not rule out failure modes for $\bar{M}$ as
$|P|$ grows: it identifies the structural quantity ($Q^*$)
through which $|P|$-dependence could enter, and shows that under
specific verification-protocol design, $Q^*$ is itself intensive in
$|P|$.  The closed form thus shifts (C7)'s epistemic content from
``$\bar{M}$ is bounded'' to ``$Q^*$ is intensive''---a more
structural claim that connects to verification-protocol design.

\section{Operational definition of \texorpdfstring{$M_{j \to j'}$}{M\_jj'}}

\subsection{Setup}

Paper 10's monitored channel structure (after (C5.MULT)) consists
of $\Kch$ channels.  Each channel $j$ has:
\begin{itemize}
\item A baseline distribution $p_0^{(j)}$ over recorded events.
\item A least-favorable adversarial alternative $p_1^{(j)}$ at
  threshold shift $\eta_j$.
\item A per-step (clipped) log-likelihood-ratio function
  $\ell^{(j)}: A \to [-\Bclip, \Bclip]$ where $A$ is the action
  space (the set of recordable event types in the verification
  ledger).  $\Bclip$ is the (C5.HOEFF) clip radius, deployment-class.
\end{itemize}

For an adversarial action distribution $q \in \Delta(A)$ (a
probability measure on $A$), define the per-step expected drift in
channel $j$:
\begin{equation}
\label{eq:drift-def}
  \mu_j(q) \;:=\; \mathbb{E}_{a \sim q}[\ell^{(j)}(a)]
  \;=\; \sum_{a \in A} q(a) \cdot \ell^{(j)}(a).
\end{equation}

The baseline-adversary $q_0$ corresponds to $p_0^{(j)}$ for all
$j$; under $q_0$, $\mu_j(q_0) = \mathbb{E}_{p_0^{(j)}}[\ell^{(j)}]$
which is approximately zero (the LLR has near-zero expectation
under the baseline).

\subsection{Cross-channel coupling magnitude}

The empirical interpretation of Audit~4 is:
``how much does a unit perturbation in channel $j$ propagate to
channel $j'$ within the cascade-time window
$\tau_{\mathrm{meta}}$.''  We formalize this as the per-step
cross-channel sensitivity:

\begin{definition}[Per-channel coupling magnitude]
\label{def:M-formal}
For ordered channel pair $(j, j')$ with $j \neq j'$, the
\emph{per-step coupling magnitude} is
\begin{equation}
\label{eq:M-def}
  M_{j \to j'}(\mathcal{D})
  \;:=\; \sup_{q \in \Aadv \cup \Aadv^{\mathrm{env}}}
  \big| \mu_{j'}(q) - \mu_{j'}(q_0) \big|,
\end{equation}
where the supremum is over the bounded adversarial class
admissible during the detection window.

The \emph{deployment coupling magnitude} is
$\bar{M}(\mathcal{D}) := \max_{j \neq j'} M_{j \to j'}(\mathcal{D})$.
\end{definition}

\begin{remark}[Why this captures Audit~4's intent]
\label{rem:M-intent}
Audit~4 measures $M$ as the response of channel $j'$ to a
perturbation in channel $j$.  Definition~\ref{def:M-formal} bounds
the response of channel $j'$ to \emph{any} adversarial choice that
deviates from baseline; this is the worst case over the
adversary's class, regardless of which channel was the
adversary's primary target.  If the adversary perturbs channel $j$
specifically, the response in $j'$ is at most
$M_{j \to j'}$; this matches Audit~4's operational reading.
\end{remark}

\section{Closed-form bound}

\subsection{The shared-action probability bound}

The key structural quantity is the maximum probability mass any
admissible adversary can place on actions that affect more than
one channel simultaneously.

\begin{definition}[Shared-action set and probability bound]
\label{def:Q-star}
For each ordered pair $(j, j')$ with $j \neq j'$, the
\emph{shared-action set} is
\begin{equation}
  A_{j, j'} \;:=\; \{ a \in A : \ell^{(j)}(a) \neq 0 \text{ and }
  \ell^{(j')}(a) \neq 0 \}.
\end{equation}

The \emph{shared-action probability bound} is
\begin{equation}
\label{eq:Q-star}
  Q^*(\mathcal{D})
  \;:=\; \max_{j \neq j'}\;
  \sup_{q \in \Aadv \cup \Aadv^{\mathrm{env}}}\;
  q(A_{j, j'})
  \;=\; \max_{j \neq j'}\; \sup_q \sum_{a \in A_{j, j'}} q(a).
\end{equation}

$Q^*(\mathcal{D}) \in [0, 1]$ is a deployment-class quantity
determined by:
\begin{enumerate}
\item The channel partition's structural overlap (i.e., how many
  actions affect more than one channel's LLR), which is fixed by
  the verification ledger's event-classification policy under
  (C5.MULT); and
\item The adversarial class $\Aadv \cup \Aadv^{\mathrm{env}}$,
  which is bounded by deployment-class parameters (rate limits,
  action bounds, channel-coverage restrictions).
\end{enumerate}
\end{definition}

\begin{remark}[Disjoint channels: $Q^* = 0$]
\label{rem:disjoint}
If the verification protocol is designed so that each action $a$
contributes to at most one channel's LLR (i.e., the channel
partition is \emph{action-disjoint}), then $A_{j, j'} = \emptyset$
for all $j \neq j'$, hence $Q^*(\mathcal{D}) = 0$.  This is the
strongest version of (C7); the bound $\bar{M} = 0$ holds trivially.

In practice, the four channels of Lemma~5b
(cooperative-output rate, attestation share, individuation drift,
bundle/cooperative-edge composition) measure structurally
different ledger event types, so action-disjointness is a
reasonable design target.  Any channel coupling that arises does
so through verification-protocol artifacts (e.g., a single ledger
event categorized into multiple channels) rather than through
intrinsic structure.
\end{remark}

\subsection{Main lemma: $\bar{M}$ in closed form}

\begin{lemma}[Closed-form bound on co-evolution coupling]
\label{lem:c7-closed-form}
Under (C5.HOEFF) (per-step LLR clipped to $[-\Bclip, \Bclip]$) and
(C5.MULT) (bounded channel multiplicity with fixed pre-deployment
partition), the per-channel coupling magnitude satisfies
\begin{equation}
\label{eq:M-bound}
  \bar{M}(\mathcal{D}) \;\leq\; 2 \cdot Q^*(\mathcal{D}) \cdot \Bclip,
\end{equation}
where $Q^*(\mathcal{D})$ is the shared-action probability bound of
Definition~\ref{def:Q-star} and $\Bclip$ is the (C5.HOEFF) clip
radius.

In particular, both $Q^*$ and $\Bclip$ are deployment-class
constants intensive in $|P|$ over $\mathcal{D}$.  The bound
\eqref{eq:M-bound} therefore proves (C7) bounded co-evolution as a
corollary of (C5.HOEFF), (C5.MULT), and the structural property
$Q^*$ of the verification protocol's channel partition.
\end{lemma}

\begin{proof}
We bound the per-step drift response in channel $j'$ to any
adversary $q \in \Aadv \cup \Aadv^{\mathrm{env}}$.

\emph{Step 1 (decompose drift contributions by action).}  Using
\eqref{eq:drift-def}:
\[
  \mu_{j'}(q) - \mu_{j'}(q_0) \;=\; \sum_{a \in A} (q(a) - q_0(a))
  \cdot \ell^{(j')}(a).
\]

\emph{Step 2 (split actions into shared and non-shared sets for
$j'$).}  Fix any other channel $j \neq j'$.  Split $A$ into:
\begin{itemize}
\item Actions shared between $j$ and $j'$:
  $A_{j, j'} = \{a : \ell^{(j)}(a) \neq 0 \text{ and }
  \ell^{(j')}(a) \neq 0 \}$.
\item Actions affecting $j'$ but not $j$:
  $A_{j'} \setminus A_{j, j'}$.
\item Actions affecting neither: contribute zero to either drift.
\end{itemize}
The drift in $j'$ decomposes as:
\[
  \mu_{j'}(q) - \mu_{j'}(q_0)
  = \underbrace{\sum_{a \in A_{j, j'}} (q(a) - q_0(a)) \ell^{(j')}(a)}_{
    \text{cross-channel coupling contribution}}
  + \underbrace{\sum_{a \in A_{j'} \setminus A_{j, j'}}
    (q(a) - q_0(a)) \ell^{(j')}(a)}_{
    \text{$j'$-only contribution}}.
\]

\emph{Step 3 (bound the cross-channel contribution).}
By (C5.HOEFF), $|\ell^{(j')}(a)| \leq \Bclip$ for all $a$.
The cross-channel contribution is therefore bounded:
\[
  \left|\sum_{a \in A_{j, j'}} (q(a) - q_0(a)) \ell^{(j')}(a)\right|
  \;\leq\; \Bclip \cdot \sum_{a \in A_{j, j'}} |q(a) - q_0(a)|
  \;\leq\; \Bclip \cdot \big(q(A_{j, j'}) + q_0(A_{j, j'})\big)
  \;\leq\; \Bclip \cdot \big(Q^*(\mathcal{D}) + Q^*(\mathcal{D})\big)
\]
where the last inequality uses that both $q$ and $q_0$ have
shared-action mass at most $Q^*(\mathcal{D})$ by
Definition~\ref{def:Q-star}.

\emph{Step 4 (the $j'$-only contribution does not couple $j$ to
$j'$).}  Actions in $A_{j'} \setminus A_{j, j'}$ have
$\ell^{(j)}(a) = 0$; perturbing these actions changes $\mu_{j'}$
but does not change $\mu_j$.  The $j'$-only contribution is
therefore not part of the cross-channel coupling from $j$ to $j'$:
under the operational reading of $M_{j \to j'}$ as the response of
$j'$ to a perturbation of $j$, only the shared-action contribution
counts, since the $j'$-only contribution would be present even if
the adversary did not perturb channel $j$.

Formally: if $q$ and $q'$ are two adversaries that agree on
$A \setminus A_{j, j'}$ (so $\mu_j(q) = \mu_j(q')$ via channel
$j$'s LLR support), but differ on $A_{j, j'}$, then
$\mu_{j'}(q) - \mu_{j'}(q')$ measures the cross-channel coupling
exactly.  By the same Step~3 argument,
\[
  |\mu_{j'}(q) - \mu_{j'}(q')|
  \;\leq\; \Bclip \cdot 2 \cdot Q^*(\mathcal{D})
  \;=\; 2 Q^* \Bclip.
\]

\emph{Step 5 (conclude).}  Taking sup over $q$ in
\eqref{eq:M-def} and noting that the bound from Step~3 covers all
admissible adversarial choices:
\[
  M_{j \to j'}(\mathcal{D}) \;\leq\; 2 Q^*(\mathcal{D}) \cdot \Bclip.
\]
Maximizing over $j \neq j'$:
\[
  \bar{M}(\mathcal{D}) \;\leq\; 2 Q^*(\mathcal{D}) \cdot \Bclip.
\]

\emph{Step 6 (intensivity of the right-hand side).}  $\Bclip$ is a
(C5.HOEFF) deployment-class constant, $|P|$-independent by
hypothesis.  $Q^*(\mathcal{D})$ depends on:
\begin{itemize}
\item The channel partition's overlap structure under (C5.MULT)'s
  fixed-cardinality $\Kch$ partition; this is fixed before
  deployment and does not change with $|P|$.
\item The verification ledger's event-classification policy, which
  is a deployment-class operational parameter.
\item The adversarial class $\Aadv \cup \Aadv^{\mathrm{env}}$,
  which is bounded by Lemma~5b/5e under (C5.MULT) familywise
  allocation.
\end{itemize}
None of these are functions of $|P|$.  Therefore $Q^*(\mathcal{D})$
is intensive in $|P|$ over $\mathcal{D}$, and so is $\bar{M}$.
\end{proof}

\section{Discussion}

\subsection{What the closed form buys}

If Lemma~\ref{lem:c7-closed-form} holds, (C7) is no longer a
free-floating assumption: it is a derived consequence of:
\begin{enumerate}
\item (C5.HOEFF) $\Bclip$ clipping --- already required for SPRT
  applicability.
\item (C5.MULT) bounded channel multiplicity --- already required
  for SPRT applicability.
\item A new structural condition on the verification protocol's
  channel partition: $Q^*(\mathcal{D}) < 1$.  This is testable from
  the verification ledger's event-classification policy and does
  not require empirical measurement of historical drift data.
\end{enumerate}

In particular, the failure mode of (C7) ($\bar{M}$ scaling with
$|P|$) reduces to the failure mode of $Q^*$: the verification
protocol's channel partition allows new ($|P|$-induced) actions to
fall into shared-action sets at rates that grow with $|P|$.  This
is a verification-protocol design question, not a structural
property of capability scaling.

\subsection{What the closed form does not buy}

\paragraph{The shared-action set must be carefully defined.}  The
proof above defines $A_{j, j'}$ as actions where both
$\ell^{(j)}(a)$ and $\ell^{(j')}(a)$ are nonzero.  This works for
the specific four-channel structure of Lemma~5b (Poisson, Bernoulli,
two multinomial), where $\ell^{(j)}(a) = 0$ corresponds to actions
that don't shift channel $j$'s baseline.  For more complex channel
structures (e.g., continuous-valued LLRs from non-discrete
distributions), the shared-action set definition may need
refinement.

\paragraph{$Q^*$ is itself a non-trivial structural property.}
The closed form shifts the load-bearing-ness from $\bar{M}$ to
$Q^*$.  $Q^*$ is more structural than $\bar{M}$ (it depends on the
channel partition rather than on observed dynamics), and is
testable from the verification protocol's event-classification
policy.  But $Q^*$ is still a structural property of the
deployment that must be measured or bounded in some way.  The
closed form is an improvement over assuming $\bar{M}$ directly
because $Q^*$ has a more natural structural interpretation (and
$Q^* = 0$ is achievable by design when channels are
action-disjoint), but it does not reduce the problem to a fully
foundational quantity.

\paragraph{Cumulative vs.\ per-step interpretation of $M$.}  The
operational reading of Audit~4 (``unit perturbation propagation
within $\tau_{\mathrm{meta}}$'') admits both per-step and
cumulative interpretations of $M$.  Lemma~\ref{lem:c7-closed-form}
above handles the per-step case.  The cumulative case multiplies
by the maximum event count $\Nev(\tau_{\mathrm{meta}}) \leq
\lambda_{\max} \cdot \tau_{\mathrm{meta}}$:
\[
  \bar{M}^{\mathrm{cumulative}} \;\leq\; 2 Q^* \Bclip \cdot
  \lambda_{\max} \cdot \tau_{\mathrm{meta}}.
\]
All three new factors ($\lambda_{\max}$, $\tau_{\mathrm{meta}}$,
the existing $Q^*$, $\Bclip$) are deployment-class constants
intensive in $|P|$.  Whether per-step or cumulative is the right
interpretation depends on how Microfoundation's
$K_{\mathrm{coev}}$ structural constant aggregates per-step
drift; this is a Microfoundation-internal question that the
closed form here does not resolve.

\subsection{Connection to (C11) capability-scaling caveat}

The vulnerability of $\bar{M}$ to capability scaling is
structurally parallel to (C11)'s capability-scaling caveat: a more
capable agent may find subtler exploitation pathways within fixed
channels.  The closed form suggests that under (C5.HOEFF) clipping
+ (C5.MULT) channel partition + bounded $Q^*$, this failure mode
is contained: subtler exploitation cannot increase $Q^*$ as long
as the verification protocol's classification policy is fixed.

A parallel closed-form derivation may apply to $\rhogap$ in (C11):
the per-step gap-growth rate is bounded by per-action LLR magnitude
$\Bclip$ times an action-coverage factor analogous to $Q^*$.  This
suggests that (C7) and (C11) can be derived together from a single
structural property of the verification protocol --- a worthwhile
direction for follow-up work that this draft does not attempt.

\section{Open questions}

\begin{enumerate}
\item \textbf{Is the shared-action probability bound $Q^*$ actually
intensive in $|P|$?}  Under (C5.MULT)'s fixed channel partition,
new actions added at higher $|P|$ are mapped into existing
channels.  But the \emph{distribution} over actions (i.e., $q$)
under an admissible adversary may concentrate more mass on shared
actions as the action space grows.  The intensivity claim depends
on a stability property of the channel partition under
$|P|$-growth that is not captured by (C5.MULT) alone.

\item \textbf{Does the per-step formulation match what
Microfoundation's $K_{\mathrm{coev}}$ actually requires?}
Microfoundation's composition uses
$(\epsnonresCoev)_+ \leq K_{\mathrm{coev}} \cdot \bar{M}$ where
$K_{\mathrm{coev}}$ is the structural co-evolution constant.  If
$K_{\mathrm{coev}}$ aggregates per-step drift over the
intervention window, the cumulative interpretation of $\bar{M}$ is
correct.  If $K_{\mathrm{coev}}$ is itself a per-step quantity,
the per-step interpretation is correct.  The Microfoundation paper
should be re-read to resolve this.

\item \textbf{Does the $j'$-only contribution argument
(Step~4) hold under arbitrary adversarial classes?}  Step~4
argues that the $j'$-only contribution does not couple $j$ to
$j'$.  This is correct for the specific Lipschitz / Audit-4
operational interpretation.  But under adversarial classes that
permit correlated perturbations across channels, the
$j'$-only contribution may track $\mu_j$ via a hidden common cause
not visible at the action level.  The proof needs to either rule
out such hidden-common-cause coupling or extend $A_{j, j'}$ to
include actions that affect $j$ via hidden mechanisms.

\item \textbf{What is the empirical value of $Q^*$ for the
canonical tripartite deployment?}  The four channels of Lemma~5b
measure substrate-distinct event types.  Reasoning suggests
$Q^*$ should be small (channels are nearly action-disjoint by
construction), but a specific deployment instance would need to
audit its channel partition for shared events.  This is
operational tooling work analogous to the existing Audit~4 but
calibrated against the new closed-form bound.
\end{enumerate}

\section{Status}

This draft is an attempt at a closed-form derivation of
$\bar{M}$.  It makes the load-bearing structural quantity $Q^*$
explicit, which is an improvement over assuming $\bar{M}$
directly.  Whether the proof is correct depends on:
\begin{itemize}
\item The Lipschitz interpretation of $M$ being the right
  operational reading (Audit~4 admits this reading; verification
  needed).
\item The $j'$-only argument (Step~4) being valid under all
  admissible adversarial classes.
\item The intensivity of $Q^*$ following from (C5.MULT) alone
  (not requiring an additional condition on channel-partition
  stability under $|P|$-growth).
\end{itemize}

The draft is suitable for codex review; codex should focus on the
proof's correctness, the cleanness of the $Q^*$ definition, and
whether the closed form actually delivers what (C7) requires.

\end{document}
